\section{模型评价与结论}

\subsection{四问结果与模型特点}

四问共用储能状态方程与结算定义，由确定性LP逐步扩展为残差场景采购、滚动候选筛选和价格--净负荷联合场景。最终日前两组保留母模型，滚动两组在零时与日内均启用筛选，结果见表\ref{tab:annual_summary}。

\begin{table}[htbp]\centering\small
\caption{最终分问配置的334天开发期闭环评价结果}
\label{tab:annual_summary}
\begin{tabularx}{\linewidth}{LRRRR}\toprule
指标 & 问题2 & 问题3 & 问题4日前版 & 问题4滚动版\\\midrule
计划费用/万元 & 1350.80 & 1350.31 & 1428.01 & 1426.52\\
调整费用/万元 & 0.00 & 24.87 & 0.00 & 28.31\\
紧急费用/万元 & 50.48 & 11.97 & 51.40 & 12.19\\
实际总费用/万元 & 1401.28 & 1387.15 & 1479.41 & 1467.02\\
紧急电量/kWh & 85017.69 & 21029.90 & 80544.00 & 21293.41\\
紧急购电天数 & 113 & 119 & 106 & 116\\
紧急购电连续时段数 & 205 & 204 & 195 & 193\\
日总费用CVaR95/元 & 70731.74 & 65552.76 & 83044.31 & 77575.87\\
最高日费用/元 & 89306.82 & 67941.61 & 108080.37 & 81797.18\\
2月1日期初库存/kWh & 8875.19 & 8814.60 & 9251.54 & 9030.40\\
12月31日期末库存/kWh & 6240.32 & 6240.32 & 6269.86 & 6274.99\\
\bottomrule\end{tabularx}
\par\smallskip\parbox{\linewidth}{\footnotesize 问题2与问题4日前版沿用母模型；问题3与问题4滚动版在零时及日内均启用候选筛选。预测器与实时执行器沿用原冻结配置，筛选适用问题与时点依据开发期比较确定。实际费用不含候选评分的终端偏离项；CVaR95为实际日总账单的经验尾部均值。}
\end{table}


问题一费用35126.95元，相对无储能节费26.90\%。第三问和第四问滚动版经筛选后，分别较母模型节省104339.41元和98994.66元，降幅为0.747\%和0.670\%；同一问题四种配置末库存一致，全时点筛选的账单最低。其收益来自计划费用节省超过新增调整费与紧急费，紧急购电量分别增至21029.90和21293.41 kWh。日前两组筛选分别增费0.161\%和0.428\%，因而继续采用母模型。

主结果逐段满足供需平衡、库存界、功率限额、充放电互斥和日际连续性，未来信息扰动检查通过。正文图表与结果附件由对应执行数组统一生成；信息来源、执行器和退款对照保留各自母模型配置，为机制解释提供依据。

\subsection{适用边界与可继续改进之处}
\label{sec:limitations}

\textbf{评价范围与选型。}结果覆盖2025年2--12月334天，预测器和实时执行器采用一月冻结配置，筛选适用问题及启用时点依据开发期费用确定。因果信息检查验证决策使用当时可得信息；独立测试与跨年泛化仍需另设未参与选型的数据。全时点筛选的优势限于本次四种配置，且伴随紧急补购及部分尾部费用增加。

\textbf{模型与回放近似。}母模型共享储能参考，候选评估则以固定执行器回放至午夜，未来预报发布与合同修订的机会价值尚未显式进入评分；内层跨日候选采用确定性续接。它们属于受限策略及有限候选评估，与完整多阶段随机最优控制仍有距离。扩大候选集可使同情景、同初态、同评分下的最小估计费用不增，未来账单仍受预测分布变化和后续状态影响。

\textbf{比较与库存口径。}同一消融内保持共同初态并独立连续运行，同时报告末库存；跨问题比较包含预测来源、初态和策略设置的共同差异。信息来源、执行器及退款实验各自回答相应机制问题，其费用变化为整个闭环运行的响应。终端偏离项是库存价值近似，仅参与优化评分；实际账单按原计划全额、调整及紧急补购三项结算。

\textbf{统计解释。}28日历史窗口提供有限经验分布，场景压缩和风险权重会共同改变采购保守程度。图中的P10--P90带表示经验分位，相关矩阵描述样本关联；置信水平、时序显著性及因果效应需额外统计设计。回放诊断的剩余时域相互重叠，费用差按决策点平均；模型中的紧急费CVaR90与结果表的实际日账单CVaR95分别对应不同损失对象。

\textbf{物理范围与后续验证。}本文采用10分钟平均功率解释、免费弃余和可按需获得的外网补购，聚焦题目给定的容量、功率及效率约束。网络容量、寿命退化、额外交易规则和更完整的跨日价值可在相应数据具备后扩展。下一步优先在未参与选型的时期检验现有分问策略，并同时评估账单、紧急补购和库存表现。
